\subsection{LLM-Based Specification, Contract, and Formal-Spec Inference}
\label{subsec:rw4}

The proposition that executable contracts can serve as the durable, machine-checkable artefact of a software system long predates the present interest in large language models (LLMs). Design by contract~\cite{meyer1992dbc} established preconditions, postconditions, and invariants as first-class obligations that bind callers and callees, while dynamic invariant detection, exemplified by Daikon~\cite{ernst2007daikon}, demonstrated that likely specifications could be \emph{recovered} from program behaviour rather than authored by hand. These two ideas---the contract as a verifiable interface and the specification as an inferable artefact---anchor the present article, but classical inference is confined to properties witnessed in observed executions and offers no bridge from the informal human intent that motivates a specification in the first place. Bridging that gap is precisely what the recent wave of LLM-based specification work attempts, and it is against this body of work that the contribution of this article must be positioned.

A first thread translates unstructured natural language directly into formal artefacts. The nl2spec framework~\cite{cosler2023nl2spec} maps English requirements to temporal logic and, recognising the inherent ambiguity of natural-language requirements, lets users iteratively correct sub-translations rather than redraft an entire formula. Closer to the contract setting, Endres et al.~\cite{endres2024nl2postcond} study \emph{nl2postcond}, prompting LLMs to convert code-adjacent natural language into formal-method postconditions, and report that the generated postconditions both correctly capture intent on a large fraction of problems and catch real historical bugs from Defects4J---direct evidence that an intent-derived oracle can disagree productively with an implementation. This last observation is foundational to the brownfield characterisation phase developed in this article, yet nl2postcond stops at producing a candidate assertion; it neither ratifies that assertion with the human nor guards against the synthesizing model silently weakening the intent it was asked to formalise.

A second, larger thread targets the specifications needed to make automated verification succeed, particularly the loop invariants that are the classical bottleneck of deductive verification. Pei et al.~\cite{pei2023invariants} show that fine-tuned models can predict invariants statically at a quality comparable to dynamic analysis, and a sequence of systems---Loopy~\cite{kamath2023loopy}, Lemur~\cite{wu2024lemur}, and AutoSpec~\cite{wen2024autospec}---couple an LLM proposer to a sound checker so that only machine-validated invariants survive. Lemur, in particular, formalizes the LLM-plus-reasoner loop as a sound calculus, while AutoSpec drives synthesis with static analysis and per-round verification to avoid error accumulation. SpecGen~\cite{ma2025specgen} extends this conversational, feedback-driven scheme to full JML method specifications for Java, falling back to mutation-based search when the model's output fails the verifier. The defining discipline of this thread is that the verifier, not the model, is the arbiter of soundness---a discipline this article adopts wholesale by reusing OpenJML and Z3 as a trusted oracle the synthesizing agent never grades.

A third thread closes the loop between specification, code, and proof. Baldur~\cite{first2023baldur} generates and repairs whole proofs with LLMs, feeding failed attempts and their error messages back into generation. Clover~\cite{sun2024clover} reduces correctness checking to mutual consistency among code, documentation, and formal annotations, using a deductive verifier together with an LLM, and the verified-Dafny synthesis of Misu et al.~\cite{misu2024dafny} generates specification and implementation together so the code can be proved against its own spec. These systems realise a counterexample-guided generate-and-check loop and confirm its feasibility, but they share a structural limitation that this article makes its central concern: the specification and the code it grades typically originate from the \emph{same} model, so the verification is, in the terminology developed here, self-grading. Where the specification is itself the output of the model under test---as in nl2postcond, Clover, and Dafny synthesis---a model that misreads the intent will produce an oracle that excuses its own implementation, and consistency among model-authored artefacts is not correctness relative to intent.

The work surveyed here thus establishes every component this article builds upon---inference of draft specifications, NL-to-formal translation, verifier-checked invariant synthesis, and the CEGIS-shaped synthesis loop---yet leaves four gaps that the proposed pipeline addresses jointly. First, none separates the authority that writes the contract from the agent that writes the code; this article makes \emph{independent authority} an explicit, provenance-auditable constraint. Second, none confronts the \emph{ratification gap}: an inferred or NL-derived contract is treated as ground truth even though the human who supplied the intent cannot read the formal text, whereas this article ratifies \emph{behaviour} through concrete pinned examples and surfaces only the boundary decisions an agent had to resolve. Third, the surveyed systems are overwhelmingly greenfield, generating fresh code against fresh specs, whereas real software is brownfield; this article recovers contracts by triangulating three disagreeing evidence streams---code, issue tickets, and documentation---and treats the disagreements as candidate bugs rather than errors to suppress. Fourth, prior work verifies a method in isolation, whereas this article promotes the compositional refinement of our prior work~\cite{edmonds_article3} to a lifecycle capability, recomputing caller obligations when a callee contract changes so that internal refactors and dependency-version bumps become the same machine-checkable blast-radius analysis~\cite{edmonds_article4}. The novelty lies not in any single inference technique but in assembling these into a verification-centric reference architecture in which a provenance- and assurance-tagged formal contract, never authored by the implementer, serves simultaneously as the agent's input and its oracle.
